\section{The interpolation operator and its normalization}\label{sec:preliminaries}
Put $Q=[-1,1]^2$ and let $\Pi_n$ be the real polynomials of total degree
at most $n$.

\noindent\textbf{Existence and uniqueness as a linear algebra question.}
A finite set of points is called \emph{unisolvent for $\Pi_n$} when
every assignment of one real value to each point is realized by exactly
one polynomial in $\Pi_n$. Equivalently, the evaluation map
\[
 \mathcal E:\Pi_n\longrightarrow\mathbb R^{\mathcal P_n},
 \qquad \mathcal E(p)=(p(B))_{B\in\mathcal P_n},
\]
is a bijection. The notation on the right means a real vector whose
coordinates are indexed by the nodes. Here the domain and target have
the same finite dimension. Consequently it suffices to prove that a
polynomial vanishing at every node must be the zero polynomial. Indeed,
that condition makes the linear map injective; a square injective
linear map is also surjective.

The standard unit vector associated with a node $B$ has entry one at
$B$ and entry zero at every other node. The cardinal polynomial $L_B$
is its inverse image under the evaluation map. Thus existence of all
the cardinal polynomials,
uniqueness of interpolation and invertibility of the evaluation matrix
are equivalent descriptions of the same fact. Lemma~\ref{lem:cardinal-kernel}
will establish that fact by computing a weighted sum of products of
basis values. The weights are positive, which will turn a vanishing
vector of values into a vanishing vector of coefficients.

The cardinal polynomials satisfy $L_B(A)=\delta_{AB}$ for
$A,B\in\mathcal P_n$, and
\[
 I_nf(z)=\sum_{B\in\mathcal P_n} f(B)L_B(z).
\]
The triangle inequality gives $\|I_n\|\leq\max_Q\lambda_n$; equality
follows by prescribing at the nodes the signs of the cardinal values
at a maximizing point and extending these finitely many values
continuously to $Q$. This proves the operator-norm interpretation used
in \eqref{eq:lebesgue-inequality}.

Here $\delta_{AB}$ is the Kronecker delta: it is one if $A=B$ and zero
otherwise. To give the norm argument in full, fix $z\in Q$. Then
\[
 |I_nf(z)|\leq\sum_B|f(B)|\,|L_B(z)|
       \leq\|f\|_\infty\lambda_n(z).
\]
Taking the maximum in $z$ gives the upper bound. For the reverse bound,
choose a point $z_*$ maximizing $\lambda_n$ and prescribe
$s_B=\operatorname{sgn}L_B(z_*)$. Recall that the sign of a positive
number is one, the sign of a negative number is minus one, and the sign
of zero is zero. The finite collection of distinct nodes has a positive minimum
separation. Choose $\rho>0$ so small that the closed balls of radius
$\rho$ around distinct nodes are disjoint, and put
\[
 \mathsf{h}_B(z)=\max\{0,1-|z-B|/\rho\},\qquad
 f_*(z)=\sum_B s_B\mathsf{h}_B(z),\qquad z\in Q.
\]
The norm $|z-B|$ here is Euclidean distance. Each $\mathsf{h}_B$ is continuous,
equals one at its own node and vanishes at the others. At any location
at most one summand is nonzero, so $|f_*(z)|\leq1$. Thus $f_*\in C(Q)$,
$f_*(B)=s_B$, and
\[
 I_nf_*(z_*)=\sum_B\operatorname{sgn}L_B(z_*)L_B(z_*)
              =\lambda_n(z_*).
\]
This constructs an admissible input attaining the bound. It also shows
why allowing continuous inputs supplies every bounded vector of nodal
values needed in the operator-norm argument.

\noindent\textbf{Why cosine polynomials are a convenient basis.}
The Chebyshev identity below specifies an ordinary algebraic polynomial
in $x=\cos t$. The first examples are
\[
 T_0(x)=1,\quad T_1(x)=x,\quad T_2(x)=2x^2-1,
 \quad T_3(x)=4x^3-3x.
\]
The cosine addition formula yields
$T_{r+1}(x)=2xT_r(x)-T_{r-1}(x)$, starting with $T_0,T_1$.
It follows inductively that $T_r$ has degree $r$ with nonzero leading
coefficient. Hence $T_0,\ldots,T_n$ form a basis for one-variable
polynomials of degree at most $n$. In two variables the products
$T_p(x)T_q(y)$ with $p+q\leq n$ form a basis for $\Pi_n$:
their leading monomials are nonzero multiples of $x^py^q$, and the
remaining terms have lower total degree. Successive elimination by
degree therefore changes between this basis and the monomial basis.
We call $(p,q)$ a \emph{mode index}; it records the two cosine
frequencies after the change to angular variables.

We use the Chebyshev polynomials $T_r(\cos t)=\cos(rt)$. Write
$\theta_j=j\pi/n$, $\phi_k=k\pi/(n+1)$ and
\[
 \varepsilon_j=\begin{cases}1/2,&j=0,n,\\1,&\text{otherwise},\end{cases}
 \qquad
 \eta_k=\begin{cases}1/2,&k=0,n+1,\\1,&\text{otherwise},\end{cases}
 \qquad w_{jk}=\frac{2\varepsilon_j\eta_k}{n(n+1)}.
\]
The endpoint factors are the half weights familiar from the trapezoidal
rule. Here they give an exact finite identity. A point interior to both angular grids has
weight $2/[n(n+1)]$; a point on one side has half that weight; a retained
corner has one quarter of that interior weight. All these numbers are
strictly positive.

The factor two compensates for retaining one color of the array. More
precisely, the \emph{checkerboard scalar product} is the weighted sum
\[
 \langle f,g\rangle_{\mathcal P_n}
   :=\sum_{\substack{0\leq j\leq n,\ 0\leq k\leq n+1\\j+k\ \mathrm{even}}}
        w_{jk}f(\cos\theta_j,\cos\phi_k)
               g(\cos\theta_j,\cos\phi_k).
\]
The expression is bilinear: holding one input fixed, it is linear in
the other input. For example,
\[
 \langle af_1+bf_2,g\rangle_{\mathcal P_n}
 =a\langle f_1,g\rangle_{\mathcal P_n}
   +b\langle f_2,g\rangle_{\mathcal P_n}
\]
for all real $a,b$. It is symmetric because exchanging $f$ and $g$
leaves every summand unchanged. It is initially positive semidefinite
on polynomials:
$\langle p,p\rangle_{\mathcal P_n}$ is a sum of nonnegative terms,
and it vanishes exactly when $p$ vanishes at every node. The lemma will
show that this form is positive definite on $\Pi_n$, meaning that it is
strictly positive for every nonzero element of that space. This is the
promised route to unisolvence. For arbitrary functions on the node set,
it is already the usual weighted Euclidean scalar product of their
vectors of values.

For the kernel formula below, set $\nu_0=1$ and $\nu_r=2$ for $r>0$;
these are the mode multiplicities. Write $z=(z_1,z_2)$ and
$B=(B_1,B_2)$ for physical points. In a sum with condition $p+q\leq n$,
both indices are nonnegative integers. If
$B=(\cos\theta_j,\cos\phi_k)$ is a node, the abbreviated notation
$w_B$ means $w_{jk}$. Define
\[
 K_n(z,B)=\sum_{p+q\leq n}\nu_p\nu_q
 T_p(z_1)T_p(B_1)T_q(z_2)T_q(B_2).
\]

The word \emph{kernel} here refers to a function of two points, formed
by summing products of matching basis functions. It is a finite
polynomial sum in $z$ and $B$. A \emph{reproducing kernel} for a scalar
product has the additional property that its scalar product with a
polynomial returns that polynomial's value at the other point. We
derive that property explicitly in the proof. The displayed $K_n$ has
the familiar coefficients $1$ for a zero frequency and $2$ for a
positive frequency in each coordinate. The grid gives one endpoint
mode a different squared norm, requiring the correction in the lemma.

One can anticipate the usual coefficients from the elementary integral
\[
 \frac1\pi\int_0^\pi\cos(pt)\cos(rt)\,dt
  =\begin{cases}1,&p=r=0,\\1/2,&p=r>0,\\0,&p\ne r.\end{cases}
\]
For a nonzero integer $k$, an antiderivative gives
\[
 \int_0^\pi\cos(kt)\,dt
   =\frac{\sin(k\pi)-\sin0}{k}=0.
\]
To obtain the displayed scalar products, use
\[
 \cos(pt)\cos(rt)
   =\tfrac12\cos((p-r)t)+\tfrac12\cos((p+r)t).
\]
When $p\ne r$, both cosine frequencies are nonzero and both integrals
vanish. When $p=r>0$, the first cosine is the constant one and the
second has nonzero frequency, giving $1/2$ after division by $\pi$.
When $p=r=0$, the integrand itself is one. The factors $\nu_r$ are
the reciprocals of these continuous squared norms. Sampling on an angular grid preserves this orthogonality for
the required indices but changes the highest available frequency in
the first coordinate. Since the two grids have sizes $n$ and $n+1$,
that change occurs asymmetrically. The proof computes it using finite
sums, so it also verifies exactly which continuous normalization survives
in the sampled setting.

The phrase \emph{cardinal kernel} in the lemma means that the corrected
kernel, with the weight of its second node attached, produces the
cardinal polynomial for that node. More explicitly, write temporarily
\[
 \widetilde K_n(z,B):=K_n(z,B)-T_n(z_1)T_n(B_1).
\]
The proof will show both that $\widetilde K_n$ reproduces $\Pi_n$ in the
checkerboard scalar product and that $w_B \widetilde K_n(A,B)=\delta_{AB}$ at
nodes. The weight belongs to the sampled variable $B$, just as a
quadrature weight belongs to the variable being summed. Later we
express this same $\widetilde K_n$ using angle sums and differences. The angular
kernel $G_n$ introduced after the proof is the two-angle trigonometric
sum from which those reflected copies of $\widetilde K_n$ are assembled. Thus
cardinal values, reproduction and the angular formula are linked by
explicit finite identities.

\begin{lemma}[Cardinal kernel]\label{lem:cardinal-kernel}
The set $\mathcal P_n$ is unisolvent for $\Pi_n$, and for
$B=(\cos\theta_j,\cos\phi_k)\in\mathcal P_n$,
\begin{equation}\label{eq:cardinal-kernel}
 L_B(z)=w_{jk}\bigl[K_n(z,B)-T_n(z_1)T_n(B_1)\bigr].
\end{equation}
In particular the coefficient of the highest first-coordinate mode is
one.
\end{lemma}
\begin{proof}
We first compute a one-dimensional weighted cosine sum. We then use
the parity condition to transfer that computation to the two-dimensional
node set. Finally we translate orthogonality into the inverse of the
interpolation system. All sums in this proof are finite.

\noindent\textbf{The endpoint-weighted geometric sum.}
For $M\geq1$ let
\[
 Q_M(g)=\frac1M\sum_{j=0}^M{}'g(j\pi/M),
\]
where the prime gives weight $1/2$ to each endpoint. The symbol $g$
denotes a function of the angle $t$. Thus $Q_M(\cos rt)$ means that
$g(t)=\cos(rt)$. The total of the weights before division by $M$ is
$M$, so $Q_M(1)=1$.

Recall the finite geometric-series identity
\[
 \sum_{j=0}^M z^j=\frac{1-z^{M+1}}{1-z}\quad(z\ne1).
\]
It follows by multiplying the sum by $1-z$: all middle powers cancel.
We use elementary complex-number notation to evaluate this sum.
A complex number is written $a+ib$, where $a,b$ are real and $i^2=-1$;
its real part is $\operatorname{Re}(a+ib)=a$. We use
$e^{i\alpha}=\cos\alpha+i\sin\alpha$ for real $\alpha$. The angle
addition formulas give $e^{i\alpha}e^{i\beta}=e^{i(\alpha+\beta)}$,
so $(e^{i\alpha})^j=e^{ij\alpha}$ and its real part is
$\cos(j\alpha)$. Only these finite algebraic identities are required
here.

To evaluate the cosine sum, let $r$ be an integer with
$1\leq r\leq2M-1$, set $z=e^{ir\pi/M}$, and write
$\sigma=z^M=(-1)^r$. In this range $z\ne1$. Half weighting the two
endpoints subtracts $(1+z^M)/2$ from the full geometric sum. Thus
\[
 \sum_{j=0}^M{}'z^j
 =\frac{1-\sigma z}{1-z}-\frac{1+\sigma}2
 =\frac{(1-\sigma)(1+z)}{2(1-z)}.
\]
For even $r$, $\sigma=1$ and this expression is zero. For odd $r$,
the quotient $(1+z)/(1-z)$ is purely imaginary. Writing
$z=e^{i\alpha}$ gives
\[
 \frac{1+e^{i\alpha}}{1-e^{i\alpha}}
   =i\cot\frac\alpha2.
\]
Indeed, factoring $e^{i\alpha/2}$ out of numerator and denominator
gives
\[
 \frac{2e^{i\alpha/2}\cos(\alpha/2)}
      {-2ie^{i\alpha/2}\sin(\alpha/2)}
 =i\frac{\cos(\alpha/2)}{\sin(\alpha/2)}.
\]
Here $1/i=-i$ follows from $i^2=-1$, and the denominator is nonzero
because $e^{i\alpha}\ne1$. The cotangent is real for these real
angles, so the expression has real part zero. Taking real parts in
the weighted geometric sum gives zero in both parity cases. Since
$\operatorname{Re}z^j=\cos(rj\pi/M)$, the weighted cosine average
vanishes. For $r=2M$ every cosine value is one, so its average is one.
This explains why the upper frequency $2M$ has to be kept separate
from the preceding vanishing range. We have proved
$Q_M(\cos rt)=0$ for $1\leq r\leq2M-1$ and
$Q_M(1)=Q_M(\cos2Mt)=1$.

Using $\cos(-t)=\cos t$, the same values apply to negative integer
frequencies. For $0\leq p,r\leq M$ the product identity
\[
 \cos(pt)\cos(rt)=\tfrac12\cos((p-r)t)
                         +\tfrac12\cos((p+r)t)
\]
now computes the scalar products. If $p\ne r$, both frequencies have
absolute value between $1$ and $2M-1$, so both averages vanish. If
$p=r$ with $0<p<M$, the constant frequency contributes $1/2$ and
the frequency $2p$ contributes zero. At $p=r=0$ both contributions
are constant, and at $p=r=M$ the frequencies are $0$ and $2M$;
the answer is one in both endpoint cases. Consequently
\[
 Q_M(\cos pt\cos rt)=\delta_{pr}h_p^{(M)},\qquad
 h_p^{(M)}=\begin{cases}1,&p=0,M,\\1/2,&0<p<M.\end{cases}
\]
Thus $h_p^{(M)}$ is the squared length of the $p$th cosine mode in this
discrete scalar product. At the highest mode, the sampled values are
$\cos(Mj\pi/M)=(-1)^j$, whose square is one at every grid point.
This gives a direct explanation of $h_M^{(M)}=1$.

\noindent\textbf{What selecting the checkerboard does.}
Put $e_{pq}(x,y)=T_p(x)T_q(y)$, where $p+q\leq n$.
The \emph{Gram matrix} is the matrix whose entry indexed by $(p,q)$
and $(r,s)$ is $\langle e_{pq},e_{rs}\rangle_{\mathcal P_n}$.
A Gram matrix records all pairwise scalar products of a chosen basis.
Its diagonal entries are the squared lengths of the basis vectors, and
an off-diagonal zero expresses orthogonality. We compute each entry
with the parity restriction included.

The indicator of an even index sum satisfies
\[
 2\mathbf1_{j+k\,\mathrm{even}}=1+(-1)^j(-1)^k.
\]
An indicator is one when its subscripted condition holds and zero
otherwise. Inserting this identity into the weighted sum extends it to
the full rectangular array and splits it into two products of
one-dimensional sums. Explicitly, the entry is
\begin{align*}
 &Q_n(\cos(pt)\cos(rt))\,
       Q_{n+1}(\cos(qt)\cos(st))\\
 &\quad+
 \left(\frac1n\sum_{j=0}^n{}'(-1)^j
                    \cos(p\theta_j)\cos(r\theta_j)\right)
 \left(\frac1{n+1}\sum_{k=0}^{n+1}{}'(-1)^k
                    \cos(q\phi_k)\cos(s\phi_k)\right).
\end{align*}
The first line is the tensor product part: ``tensor product'' here
means the product of an independent sum in each coordinate. The
second line is the alternating part caused by the choice of color.

At the grid angles $t_j=j\pi/M$, the identity
\[
 \cos((M-p)t_j)
 =\cos(Mt_j)\cos(pt_j)+\sin(Mt_j)\sin(pt_j)
 =(-1)^j\cos(pt_j)
\]
holds because $\sin(j\pi)=0$. Thus multiplication by $(-1)^j$
sends mode $p$ to mode $M-p$ on the grid. The first alternating factor
is $Q_n(\cos((n-p)t)\cos(rt))$, which can be nonzero only if
$r=n-p$. The second factor can be nonzero only if $s=n+1-q$.
All these indices lie in the one-dimensional orthogonality ranges
already proved. In particular $q,s\leq n<n+1$.

For the product to survive, both necessary equalities must hold at once.
They would imply
\[
 p+r=n,\qquad q+s=n+1,\qquad p+r+q+s=2n+1.
\]
But the two mode indices are admissible, so
$(p+q)+(r+s)\leq2n$. This contradiction makes the alternating part
zero. It is exactly here that the triangular mode constraint and the
two consecutive grid sizes work together. The tensor product part is
computed by the one-dimensional identity. Therefore
\[
 \langle e_{pq},e_{rs}\rangle_{\mathcal P_n}
 =\delta_{pr}\delta_{qs}h_p^{(n)}h_q^{(n+1)}.
\]
The Gram matrix is diagonal and positive, with diagonal
$h_p^{(n)}h_q^{(n+1)}$. In particular the $(0,0)$ entry equals one,
so the nodal weights sum to one.

\noindent\textbf{Positive Gram matrix and interpolation.}
Expand a polynomial as $P=\sum_{p+q\leq n}c_{pq}e_{pq}$.
The orthogonality identity gives
\[
 \langle P,P\rangle_{\mathcal P_n}
  =\sum_{p+q\leq n}c_{pq}^2h_p^{(n)}h_q^{(n+1)}.
\]
Each factor on the right is positive. If $P$ vanishes at every node,
the left side is zero. Each coefficient must therefore vanish, proving
injectivity of the evaluation map. There are $(n+1)(n+2)/2$
checkerboard points, the same number as the modes, so the evaluation
map is also surjective. This proves unisolvence.

In matrix notation, choose any fixed ordering of the nodes and mode
indices, and let
\[
 E_{B,(p,q)}=e_{pq}(B),\qquad
 W=\operatorname{diag}(w_B),\qquad
 D=\operatorname{diag}(h_p^{(n)}h_q^{(n+1)}).
\]
Here $\operatorname{diag}$ denotes the matrix with the listed entries
on its diagonal and zeros elsewhere; $w_B$ is $w_{jk}$ at node $B$.
The transpose $E^{\mathsf T}$ exchanges row and column indices.
Matrix multiplication sums products over the shared index. Since $W$
is diagonal, the entry of the product indexed by the two modes is
\begin{align*}
 (E^{\mathsf T}WE)_{(p,q),(r,s)}
 &=\sum_{B\in\mathcal P_n}
       E_{B,(p,q)}w_BE_{B,(r,s)}\\
 &=\sum_{B\in\mathcal P_n}w_B e_{pq}(B)e_{rs}(B)
   =\langle e_{pq},e_{rs}\rangle_{\mathcal P_n}.
\end{align*}
Our orthogonality calculation makes this entry zero for distinct modes
and $h_p^{(n)}h_q^{(n+1)}$ for identical modes. These are exactly the
entries of $D$, proving $E^{\mathsf T}WE=D$. To see explicitly why
this implies independent columns, suppose a coefficient vector $c$
satisfies $Ec=0$. Then
\[
 c^{\mathsf T}Dc=(Ec)^{\mathsf T}W(Ec)=0.
\]
The left side is a sum of strictly positive diagonal factors times the
squares of the entries of $c$. Consequently all those entries vanish,
which is precisely column independence. Since $E$ is square, it is invertible. This matrix
expression also tracks the kernel normalization.

For $q\leq n$, the inverse of $h_q^{(n+1)}$ is $\nu_q$. In the
first coordinate the inverse of $h_p^{(n)}$ is $\nu_p$ for $p<n$,
while the inverse at $p=n$ is one. The only admissible mode with $p=n$
has $q=0$. Thus the inverse Gram coefficients are $\nu_p\nu_q$ except
at $(p,q)=(n,0)$, where the coefficient is one. Precisely one product
coefficient differs from the continuous normalization described above.

\noindent\textbf{Why a reproducing kernel yields the cardinal values.}
Define the kernel with these inverse coefficients by
\[
 \widetilde K_n(z,B)=\sum_{p+q\leq n}
   \frac{e_{pq}(z)e_{pq}(B)}{h_p^{(n)}h_q^{(n+1)}}
        =K_n(z,B)-T_n(z_1)T_n(B_1).
\]
For a fixed $z$, take its scalar product in the $B$ variable with
$P(B)=\sum_{r+s\leq n}c_{rs}e_{rs}(B)$. Substituting both finite
expansions and interchanging finite sums gives
\begin{align*}
 &\sum_{B\in\mathcal P_n}w_B \widetilde K_n(z,B)P(B)\\
 &\quad=\sum_{p+q\leq n}\sum_{r+s\leq n}
    \frac{e_{pq}(z)c_{rs}}{h_p^{(n)}h_q^{(n+1)}}
       \sum_{B\in\mathcal P_n}w_Be_{pq}(B)e_{rs}(B)\\
 &\quad=\sum_{p+q\leq n}\sum_{r+s\leq n}
    \frac{e_{pq}(z)c_{rs}}{h_p^{(n)}h_q^{(n+1)}}
          \langle e_{pq},e_{rs}\rangle_{\mathcal P_n}.
\end{align*}
The inner scalar product is zero unless $r=p$ and $s=q$. For those
surviving indices its value is $h_p^{(n)}h_q^{(n+1)}$, exactly
canceling the denominator. Hence the double sum reduces to the single
basis expansion and yields
\[
 \sum_{B\in\mathcal P_n}w_B \widetilde K_n(z,B)P(B)
      =\sum_{p+q\leq n}c_{pq}e_{pq}(z)=P(z).
\]
This is the reproducing property: the weighted scalar product with
$\widetilde K_n(z,\cdot)$ returns evaluation at $z$ for every $P\in\Pi_n$.

To see the nodal delta values directly, let $F$ be the matrix with
entries $F_{A,B}=w_B \widetilde K_n(A,B)$. To obtain its matrix formula, compute
one entry of the proposed product:
\begin{align*}
 (ED^{-1}E^{\mathsf T}W)_{A,B}
 &=w_B\sum_{p+q\leq n}
       \frac{e_{pq}(A)e_{pq}(B)}{h_p^{(n)}h_q^{(n+1)}}\\
 &=w_B\widetilde K_n(A,B)=F_{A,B}.
\end{align*}
The equality of all entries proves the formula for $F$. Associating
matrix factors and then using $E^{\mathsf T}WE=D$ gives
\[
 F=ED^{-1}E^{\mathsf T}W,
 \qquad FE=ED^{-1}(E^{\mathsf T}WE)=E.
\]
Multiplying $FE=E$ on the right by $E^{-1}$ gives $F$ equal to the
identity matrix. Hence $w_B \widetilde K_n(A,B)=\delta_{AB}$. For fixed $B$,
the function $w_B \widetilde K_n(z,B)$ lies in $\Pi_n$ and has precisely the
prescribed cardinal values. Uniqueness identifies it with $L_B$ and
proves \eqref{eq:cardinal-kernel}. The exceptional coefficient in $K_n$
was $\nu_n\nu_0=2$; subtracting one copy of its product leaves
coefficient one, as claimed.
\end{proof}

\noindent\textbf{A complete small example of the normalization.}
The case $n=1$ provides a short illustration of the kernel construction.
Although the main location theorem is stated from degree two onward,
the interpolation identities in this section already apply at degree
one. The nodes are
\[
 B_1=(1,1),\qquad B_2=(1,-1),\qquad B_3=(-1,0),
\]
with respective weights $1/4,1/4,1/2$. In the basis $1,x,y$, the
weighted Gram matrix is diagonal with entries $1,1,1/2$. For instance,
\[
 \langle1,x\rangle_{\mathcal P_1}
       =\tfrac14+\tfrac14-\tfrac12=0,
 \qquad
 \langle x,x\rangle_{\mathcal P_1}=1,
 \qquad
 \langle y,y\rangle_{\mathcal P_1}=\tfrac12.
\]
The distinction between the last two squared norms reflects the two
unequal grids.

For $z=(x,y)$ and $B=(\xi,\eta)$, the uncorrected sum and the
corrected sum are
\[
 K_1(z,B)=1+2x\xi+2y\eta,\qquad
 R_1(z,B)=1+x\xi+2y\eta.
\]
Multiplying the corrected expression by each nodal weight gives
\[
 L_{B_1}(x,y)=\frac{1+x+2y}{4},\qquad
 L_{B_2}(x,y)=\frac{1+x-2y}{4},\qquad
 L_{B_3}(x,y)=\frac{1-x}{2}.
\]
Substituting the three node coordinates verifies all nine cardinal
values directly. Summing the formulas gives one. The exceptional
$x\xi$ coefficient is visibly one: replacing it by two would alter
the nodal values and the interpolating polynomial. This example shows
how a single normalization coefficient affects the entire construction.

It also illustrates the absolute values in the Lebesgue function.
For real $a,b$,
\[
 |a+b|+|a-b|=2\max(|a|,|b|).
\]
To prove this, first suppose $|a|\geq|b|$. If $a\geq0$, then both
$a+b$ and $a-b$ are nonnegative and their absolute values sum to
$2a=2|a|$. If $a<0$, both are nonpositive and the sum is $-2a=2|a|$.
In the other case, $|b|\geq|a|$, exchange $a$ and $b$ in that
argument; the left side is unchanged because $|b-a|=|a-b|$, and the
answer is $2|b|$. This includes equality and zero cases.

For the first two cardinal polynomials, apply the identity with
$a=1+x$ and $b=2y$ and divide by four. On $Q$, $1+x$ and $1-x$ are
nonnegative, so the third cardinal polynomial already has nonnegative
value. Therefore
\[
 \lambda_1(x,y)
 =\max\left(\frac{1+x}{2},|y|\right)+\frac{1-x}{2}
 =\max\left(1,|y|+\frac{1-x}{2}\right).
\]
For every $(x,y)\in Q$, both $|y|\leq1$ and $(1-x)/2\leq1$,
so the second entry in this maximum is at most two. It equals two
precisely when both bounds are equalities: $|y|=1$ and $x=-1$.
The first entry is one and cannot reach two. Thus the largest value
is two, attained exactly at $(-1,1)$ and $(-1,-1)$.
All of this follows from the displayed affine polynomials. For larger
$n$ there are many more changing signs; the kernel representation
organizes those signs in a way suited to the later proofs.

\noindent\textbf{Passing from physical coordinates to angles.}
Every $z=(x,y)\in Q$ can be written uniquely as
$z=(\cos\theta,\cos\phi)$ with $\theta,\phi\in[0,\pi]$.
In these variables every $T_p(x)T_q(y)$ becomes
$\cos(p\theta)\cos(q\phi)$. A polynomial kernel of two physical
points therefore becomes a trigonometric expression involving four
angles. The \emph{angular kernel} is the finite sum in two free angle
variables $u,v$ that lets us express the four-angle kernel as an average
of reflected copies. Its definition below includes the same exceptional
mode correction already determined by the scalar product.

The same correction appears in the angular kernel
\begin{equation}\label{eq:angular-kernel}
 G_n(u,v)=\sum_{p+q\leq n}\nu_p\nu_q\cos(pu)\cos(qv)-\cos(nu).
\end{equation}
At a corner it gives the exact signed row sums; in the interior its
reflections give the cancellation comparison. Thus the two arguments
start from one normalized kernel.

To spell out the relation, write $B=(\cos\alpha,\cos\beta)$ and
$z=(\cos\theta,\cos\phi)$. The product-to-sum identity gives
\[
 \cos(p\theta)\cos(p\alpha)
    =\tfrac12\{\cos(p(\theta-\alpha))+
                         \cos(p(\theta+\alpha))\}.
\]
Applying it in both coordinates, including the corrected term, yields
\[
 \widetilde K_n(z,B)=\frac14\sum_{\sigma\in\{-1,1\}}
                    \sum_{\tau\in\{-1,1\}}
         G_n(\theta+\sigma\alpha,\phi+\tau\beta).
\]
The correction $-\cos(nu)$ in $G_n$ contributes
$-\cos(n\theta)\cos(n\alpha)$ to this average, because the sum in
$\tau$ simply supplies two equal copies. Thus it reproduces exactly
the correction in $\widetilde K_n$. Substituting node angles $\alpha=\theta_j$
and $\beta=\phi_k$ and multiplying by $w_{jk}$ recovers $L_B(z)$.

The function $G_n$ is a trigonometric polynomial: a finite linear
combination of products of sine or cosine functions, here consisting
only of cosine products. It is even in each angle and $2\pi$-periodic
in each angle. These properties justify reducing angle expressions
modulo $2\pi$ and reflecting them when convenient. Such reductions
apply to the angle arguments of $G_n$; the two different grid mesh
widths remain fixed throughout.

This representation explains why both angle sums and angle differences
appear in later formulas. They arise from multiplying a mode evaluated
at $z$ by the same mode evaluated at $B$. At a corner, $\theta$ and
$\phi$ are each either zero or $\pi$, so cosine values simplify to
signs and the two reflected contributions in a coordinate are closely
related. At an interior point their relative signs carry additional
information. For real numbers $a_1,\ldots,a_m$ the nonnegative quantity
\[
 \sum_{r=1}^m|a_r|-\left|\sum_{r=1}^m a_r\right|
\]
measures the cancellation lost if the absolute value is moved inside
the sum. Applying the triangle inequality to the four reflected copies
of $G_n$ produces precisely this kind of loss. Later comparisons retain
and estimate that loss, because its size helps distinguish a corner
from a competing point.

\noindent\textbf{Why the surviving reflections preserve interpolation.}



The two nonnode corners in Theorem~\ref{thm:main-location} have equal
Lebesgue values by symmetry. For even $n$, reflection in the first
coordinate sends $j$ to $n-j$ and preserves the checkerboard. For odd
$n$, reflection in the second coordinate sends $k$ to $n+1-k$ and
preserves it. Uniqueness of interpolation carries these symmetries to
the cardinal polynomials and hence to $\lambda_n$.

Here is the uniqueness argument explicitly. Let $S$ denote the relevant
coordinate reflection, so $S$ permutes $\mathcal P_n$ and $S^2$ is the
identity map. Composition with $S$ preserves total degree. At any node
$A$,
\[
 L_{SB}(SA)=\delta_{SA,SB}=\delta_{A,B}.
\]
Thus the polynomial $z\mapsto L_{SB}(Sz)$ has the same cardinal data
as $L_B(z)$ and belongs to the same space. Unisolvence gives
$L_{SB}(Sz)=L_B(z)$ for every $z\in Q$. Reindexing the sum of absolute
values by the node permutation $B\mapsto SB$ proves
\[
 \lambda_n(Sz)=\sum_{B\in\mathcal P_n}|L_{SB}(Sz)|
              =\sum_{B\in\mathcal P_n}|L_B(z)|=\lambda_n(z).
\]
For even $n$, the map $(j,k)\mapsto(n-j,k)$ changes $j+k$ by an even
integer and corresponds to $(x,y)\mapsto(-x,y)$. For odd $n$, the map
$(j,k)\mapsto(j,n+1-k)$ has the same parity property and corresponds
to $(x,y)\mapsto(x,-y)$. These maps exchange the two nonnode corners
listed in Theorem~\ref{thm:main-location}. This proves their equality
of value before the later strict global comparison is undertaken.

\noindent\textbf{Continuous values at quotient formulas.}
Finite cosine sums are continuous at every real angle. When a geometric
sum is written as a trigonometric quotient, its denominator can vanish
at some angles even though the original finite sum is perfectly
well-defined. A simple example is
\[
 \sum_{r=0}^{M-1}e^{irt}
    =e^{i(M-1)t/2}\frac{\sin(Mt/2)}{\sin(t/2)}
      \quad(t\notin2\pi\mathbb Z).
\]
At $t=0$ the left side equals $M$, and the quotient on the right has
that same limit. The polynomial or finite trigonometric expression
specifies the value at the exceptional angle. This convention applies
in particular when an evaluation angle meets a node angle, or when a
sum or difference of angles reaches an endpoint. Such events are called
\emph{collisions} below. They are included in every closed domain on
which a bound is asserted.


\begin{remark}
All formulas involving trigonometric quotients are interpreted by their
continuous polynomial limits at removable singularities. The maximization
and certificate domains include all nodes, edges and collisions of angular
coordinates.
\end{remark}
